\scanpage{24}
% L1, light pencil
\margin{$\deleted{\unsure{(X/Y)^p}}=\unsure{(X/Y)^p}=\infty\qquad\unsure{Z|2}$}

% L1–3, upper left
\diagram{\begin{tikzcd}[ampersand replacement=\&]
\&Y\arrow[d,"p"]\&\pi^{(n)}(Y/X/S)\arrow[d]\\
S\&X\arrow[l,"f"]\&\Delta^{(n)}_{X/S}\arrow[l,"{\unsure{g|p|D|\zeta|\delta}}^{n}_{X/S}"]
\end{tikzcd}}
% L3–5
On veut définir les ``fibrés''
\[
p^n(Y/X/S)
\]
des \deleted{éléments} ``éléments de sections d'ordre $n$ de $Y/S$ \unsure{en|sur} $X/S$''.

% L5–9, left
\diagram{\begin{tikzcd}[ampersand replacement=\&]
\&X\times_S Y\arrow[d]\\
X\arrow[r,"\Delta_X"]\&X\times_S X\arrow[d,"\operatorname{pr}_1"]\\
\&X
\end{tikzcd}}
% L6–9
\inserted{$X\times_S X$ est} muni des deux projections $\operatorname{pr}^{\deleted{\unsure{2}}}_1$, $\operatorname{pr}_2$ et de l'application diagonale $\Delta_X$. On considère $\deleted{\unsure{\delta}}\ \Delta^{(n)}_{X/S}$ le $n$-ième voisinage infinitésimal de l'immersion $\Delta_X$ comme un $X$-\deleted{\unsure{p|pré}}schéma grâce à $\operatorname{pr}_1$.

% L9–12, left
\diagram{\begin{tikzcd}[ampersand replacement=\&]M\arrow[d]\\
{\unsure{P}}^{n}_{A/\Lambda}\otimes_A M\arrow[d]\\ M
\end{tikzcd}}
% L10–13
$i^{(n)}_X$ désigne l'immersion de $\Delta^{(n)}_{X/S}$ \deleted{dans} $X\times_S X$. On considère alors le
\[
\deleted{p^n(Y/X/S)=(\operatorname{pr}_2 i^{(n)}_X)^\ast}
\]
$\Delta^{(n)}_{X/S}$-schéma $(\operatorname{pr}_2 i^{(n)}_X)^\ast(Y/X)\inserted{=\Delta^n_{X/S}\times_{X\times_S X}X\times_S Y}=\deleted{\unsure{p}}\pi^{(n)}(Y/X/S)$

% L14–17, diagram at left
\diagram{\begin{tikzcd}[ampersand replacement=\&,column sep=small,row sep=small,execute at end picture={\draw ([xshift=-4pt,yshift=4pt]\tikzcdmatrixname-1-3.north west) rectangle ([xshift=4pt,yshift=-5pt]\tikzcdmatrixname-2-4.south east);}]
\vphantom X\&X\times_S Y\arrow[l]\arrow[d]\&p^n(Y/X/S)\arrow[l]\arrow[d]\&Y\arrow[l,"{\varepsilon}^{n}(Y/X/S)"']\arrow[d]\\
X\&X\times_S X\arrow[l,"\operatorname{pr}_1"']\&\Delta^{(n)}_{X/S}\arrow[l,"i^{(n)}_{X/S}"']\&X\arrow[l]\\
X'\arrow[u,"\varphi"]\&X'\times_S X\arrow[u]\arrow[l,"\unsure{\mathrm{id}|i\Lambda}"']\&\&X'\arrow[u]\arrow[ll,"(\Gamma_{\varphi})^{(n)}"']
\end{tikzcd}}
% L14–16
On a donc $X\times_{\Delta^{(n)}_{X/S}}\pi^n(Y/X/S)=Y$, d'où $Y\longrightarrow\pi^n(Y/X/S)$ compatible avec
\[
X\xrightarrow{\Delta_X}\Delta^{(n)}_{X/S}
\]
% L17–19
Il y a aussi un \deleted{le} morphisme composé
\[
d^{(n)}_{Y/X/S}:\pi^n(Y/X/S)\longrightarrow X\times_S Y\longrightarrow Y
\]
et on a commutativité de
% L20–21
\diagram{\begin{tikzcd}[ampersand replacement=\&,execute at end picture={\draw ([xshift=-5pt,yshift=5pt]\tikzcdmatrixname-1-2.north west) rectangle ([xshift=5pt,yshift=-5pt]\tikzcdmatrixname-2-3.south east);}]
Y\arrow[r,"\varepsilon^n(Y/X/S)"]\arrow[d]\&\pi^n(Y/X/S)\arrow[d]\arrow[r,"d^n(Y/X/S)"]\&Y\arrow[d]\\
X\arrow[r,"\Delta_{X/S}"]\&\Delta^{(n)}_{X/S}\arrow[r,"d^{(n)}_{X/S}"]\&\deleted{\unsure{\Delta^n_{X/S}}}X
\end{tikzcd}}
% archive sheet number: 23

